\section{Fourier transform}

\subsection{From a series to an integral}

For a nonperiodic function on the real line, regard its Fourier series as the limit $\ell\to\infty$. With $\omega_n=n\pi/\ell$ and $\Delta\omega=\pi/\ell$, the discrete frequencies become continuous. The original source keeps the limiting idea explicit:
\begin{equation}
f(x)=\lim_{\ell\to\infty}\sum_{n=-\infty}^{\infty} c_n e^{\mathrm{i}\omega_n x},
\qquad
c_n=\frac{1}{2\ell}\int_{-\ell}^{\ell}f(x)e^{-\mathrm{i}\omega_n x}\,dx.
\end{equation}
As $\ell\to\infty$, the sum turns into an integral, and the coefficient becomes frequency density. This is the essence of the Fourier integral. Using the symmetric convention,
\begin{align}
F(\omega)&=\frac{1}{\sqrt{2\pi}}\int_{-\infty}^{\infty}f(x)e^{-\mathrm{i}\omega x}\,dx,\\
f(x)&=\frac{1}{\sqrt{2\pi}}\int_{-\infty}^{\infty}F(\omega)e^{\mathrm{i}\omega x}\,d\omega.
\end{align}
The Fourier integral theorem holds when $f$ satisfies the Dirichlet conditions on every finite interval and is absolutely integrable on the real line. At a discontinuity it returns the mean of the two limiting values.

Equivalently,
\begin{equation}
f(x)=\int_0^\infty A(\omega)\cos(\omega x)\,d\omega+
\int_0^\infty B(\omega)\sin(\omega x)\,d\omega,
\end{equation}
where $A(\omega)=\pi^{-1}\int f(x)\cos(\omega x)\,dx$ and $B(\omega)=\pi^{-1}\int f(x)\sin(\omega x)\,dx$. The representation $C(\omega)\cos[\omega x-\varphi(\omega)]$ defines the amplitude spectrum $C=(A^2+B^2)^{1/2}$ and phase spectrum $\varphi=\arctan(B/A)$.

\begin{example}
For $f(x)=e^{-\alpha|x|}$ with $\alpha>0$, we split the integral at $x=0$ and obtain
\begin{align}
F(\omega) &=\sqrt{\frac{1}{2\pi}}\left[\int_{-\infty}^{0}e^{\alpha x}e^{-\mathrm{i}\omega x}\,dx+
\int_0^{\infty}e^{-\alpha x}e^{-\mathrm{i}\omega x}\,dx\right]\nonumber\\
&=\sqrt{\frac{1}{2\pi}}\left(\frac{1}{\alpha+\mathrm{i}\omega}+\frac{1}{\alpha-\mathrm{i}\omega}\right)
onumber\\
&=\sqrt{\frac{1}{2\pi}}\frac{2\alpha}{\alpha^2+\omega^2}.
\end{align}
The transform is real and even because the original function is even. For the delta distribution, $\mathcal{F}[\delta(x)]=1/\sqrt{2\pi}$. For $f(x)=e^{-\alpha x^2}$,
\begin{equation}
F(\omega)=\frac{1}{\sqrt{2\alpha}}e^{-\omega^2/(4\alpha)}.
\end{equation}
Thus a Gaussian transforms into a Gaussian, and greater localization in $x$ corresponds to greater spread in frequency.
\end{example}

\subsection{Basic properties}

If $\mathcal{F}[f]=F$, then
\begin{align}
\mathcal{F}[f']&=\mathrm{i}\omega F(\omega), &
\mathcal{F}\left[\int^x f(s)\,ds\right]&=\frac{F(\omega)}{\mathrm{i}\omega},\\
\mathcal{F}[f(ax)]&=\frac{1}{|a|}F\left(\frac{\omega}{a}\right), &
\mathcal{F}[f(x-x_0)]&=e^{-\mathrm{i}\omega x_0}F(\omega),\\
\mathcal{F}[e^{\mathrm{i}\omega_0x}f(x)]&=F(\omega-\omega_0).
\end{align}
The derivative theorem is especially important: differentiation in physical space corresponds to multiplication by $\mathrm{i}\omega$ in frequency space. The same idea converts many linear ODEs and PDEs into algebraic problems in Fourier space. For convolution $(f_1*f_2)(x)=\int_{-\infty}^{\infty}f_1(s)f_2(x-s)\,ds$, the symmetric normalization gives $\mathcal{F}[f_1*f_2]=\sqrt{2\pi}F_1F_2$.

\subsection{Higher-dimensional transforms}

For $f(\mathbf{r})$ on $\mathbb{R}^d$,
\begin{align}
F(\mathbf{k})&=\frac{1}{(2\pi)^{d/2}}\int_{\mathbb{R}^d}f(\mathbf{r})e^{-\mathrm{i}\mathbf{k}\cdot\mathbf{r}}\,d^d\mathbf{r},\\
f(\mathbf{r})&=\frac{1}{(2\pi)^{d/2}}\int_{\mathbb{R}^d}F(\mathbf{k})e^{\mathrm{i}\mathbf{k}\cdot\mathbf{r}}\,d^d\mathbf{k}.
\end{align}
For a rectangular aperture $f(x,y)=A$ on $|x|\leq X$, $|y|\leq Y$ and zero otherwise, its two-dimensional transform is proportional to
\begin{equation}
AXY\frac{\sin(k_1X)}{k_1X}\frac{\sin(k_2Y)}{k_2Y}.
\end{equation}
This sinc pattern is the Fraunhofer diffraction pattern of a rectangular opening.